% ============================================================================
%  lessons/lesson02.tex — درسِ دوم: اَلْوَجْهُ النَّافِعُ وَ الْوَجْهُ الْمُضِرُّ
%  محتوا: متنِ کتاب (داستانِ نوبل) + ترجمه + معجم + قواعد (الحال و واوِ حالیه)
%         + خودآزمایی و تمرین‌های کتاب + پاسخ‌نامهٔ تشریحی
% ============================================================================

\Lesson
  {اَلدَّرْسُ الثَّانِي}
  {اَلْوَجْهُ النَّافِعُ وَ الْوَجْهُ الْمُضِرُّ}
  {أَعْقَلُ النَّاسِ أَنْظَرُهُمْ فِي الْعَوَاقِبِ.}
  {خردمندترِ مردم، تیزبین‌ترِ ایشان در فرجامِ [کارها] است.}
  {اَلْإِمَامُ عَلِيٌّ (ع)}

\JGoals{مفهومِ \emph{اَلْحَال} (قیدِ حالت) و تفاوتِ آن با صفت \JGoalSep جملهٔ
حالیّه با واوِ حالیّه \JGoalSep ترجمهٔ متنِ داستانی \JGoalSep تحلیلِ صرفی و
اعرابِ واژه‌ها}

% ============================================================================
%  ۱) متنِ درس
% ============================================================================
\Jsec{اَلنَّصُّ \,\textbar\, متنِ درس}

\begin{JAyat}
\Ar ﴿إِنَّ الَّذِينَ آمَنُوا وَ عَمِلُوا الصَّالِحَاتِ إِنَّا لا نُضِيعُ أَجْرَ
مَنْ أَحْسَنَ عَمَلًا﴾ \hfill اَلْكَهْف: ٣

بی‌گمان کسانی که ایمان آوردند و کارهای شایسته کردند — [به‌راستی] ما پاداشِ
کسی را که نیکوکاری کرده است تباه نمی‌کنیم.
\end{JAyat}

% ------------------------------------------------------ بندِ اوّل -----------
\begin{JNass}[الصَّبِيُّ الْمُبْدِعُ \,\text| پسرِ نابغه]
في عامِ أَلْفٍ وَ ثَمَانِمِائَةٍ وَ ثَلَاثَةٍ وَ ثَلَاثِينَ وُلِدَ فِي مَمْلَكَةِ
«السُّوَيْد» صَبِيٌّ سُمِّيَ «أَلْفِرِد نُوبِل». كَانَ وَالِدُهُ قَدْ أَقَامَ
مَصْنَعًا لِصِنَاعَةِ مَادَّةِ «النِّيتْرُوغْلِيسِرِين» السَّائِلِ السَّرِيعِ
الاِنْفِجَارِ، وَلَوْ بِالْحَرَارَةِ الْقَلِيلَةِ.\FNmark{١}

اِهْتَمَّ أَلْفِرِد مُنْذُ صِغَرِهِ بِهٰذِهِ الْمَادَّةِ، وَ عَمِلَ عَلَى
تَطْوِيرِهَا مُجِدًّا، لِيَمْنَعَ انْفِجَارَهَا. بَنَى مُخْتَبَرًا صَغِيرًا
لِيُجْرِيَ فِيهِ تَجَارِبَهُ، وَلٰكِنْ مَعَ الْأَسَفِ انْفَجَرَ الْمُخْتَبَرُ
وَ انْهَدَمَ عَلَى رَأْسِ أَخِيهِ الْأَصْغَرِ وَ قَتَلَهُ. هٰذِهِ الْحَادِثَةُ
لَمْ تُضْعِفْ عَزْمَهُ، فَقَدْ وَاصَلَ عَمَلَهُ دَؤُوبًا، حَتَّى اسْتَطَاعَ
أَنْ يَخْتَرِعَ مَادَّةَ «الدَّينَامِيت» الَّتِي لا تَنْفَجِرُ إِلَّا
بِإِرَادَةِ الْإِنْسَانِ. بَعْدَ أَنِ اخْتَرَعَ الدَّينَامِيتَ، أَقْبَلَ عَلَى
شِرَائِهِ رُؤَسَاءُ شَرِكَاتِ الْبِنَاءِ وَ الْمَنَاجِمِ وَ الْقُوَّاتُ
الْمُسَلَّحَةُ وَ هُمْ مُشْتَاقُونَ لِاسْتِخْدَامِهِ، فَانْتَشَرَ الدَّينَامِيتُ
فِي جَمِيعِ أَنْحَاءِ الْعَالَمِ.

\FNfoot{١}{كَانَ وَالِدُهُ قَدْ أَقَامَ مَصْنَعًا: پدرش کارخانه‌ای ساخته بود.}
\end{JNass}

% ------------------------------------------------------ بندِ دوم -----------
\begin{JNass}[الثَّرْوَةُ وَ الْإِنْجَازُ \,\text| ثروت و دستاورد]
قَامَ أَلْفِرِدُ بِإِنْشَاءِ عَشَرَاتِ الْمَصَانِعِ وَ الْمَعَامِلِ فِي عِشْرِينَ
دَوْلَةً، وَ كَسَبَ مِنْ ذٰلِكَ ثَرْوَةً كَبِيرَةً جِدًّا حَتَّى أَصْبَحَ مِنْ
أَغْنَى أَغْنِيَاءِ الْعَالَمِ.

فَقَدِ اسْتَفَادَ الْإِنْسَانُ مِنْ هٰذِهِ الْمَادَّةِ، وَ سَهَّلَ أَعْمَالَهُ
الصَّعْبَةَ فِي حَفْرِ الْأَنْفَاقِ وَ شَقِّ الْقَنَوَاتِ وَ إِنْشَاءِ
الطُّرُقِ وَ حَفْرِ الْمَنَاجِمِ وَ تَحْوِيلِ الْجِبَالِ وَ التِّلَالِ إِلَى
سُهُولٍ صَالِحَةٍ لِلزِّرَاعَةِ.

وَ مِنَ الْأَعْمَالِ الْعَظِيمَةِ الَّتِي تَمَّتْ بِوَاسِطَةِ هٰذِهِ الْمَادَّةِ
تَفْجِيرُ الْأَرْضِ فِي قَنَاةِ «بَنَمَا» بِمِقْدَارٍ مِنَ الدَّينَامِيتِ بَلَغَ
أَرْبَعِينَ طُنًّا.

وَ بَعْدَ أَنِ اخْتَرَعَ نُوبِلُ الدَّينَامِيتَ، اِزْدَادَتِ الْحُرُوبُ وَ
كَثُرَتْ أَدَوَاتُ الْقَتْلِ وَ التَّخْرِيبِ بِهٰذِهِ الْمَادَّةِ، وَ إِنْ كَانَ
غَرَضُهُ مِنِ اخْتِرَاعِهِ مُسَاعَدَةَ الْإِنْسَانِ فِي مَجَالِ الْإِعْمَارِ وَ
الْبِنَاءِ.
\end{JNass}

% ------------------------------------------------------ بندِ سوم -----------
\begin{JNass}[جَائِزَةُ نُوبِل \,\text| جایزهٔ نوبل]
نَشَرَتْ إِحْدَى الصُّحُفِ الْفَرَنْسِيَّةِ عِنْدَ مَوْتِ أَخِيهِ الْآخَرِ
عُنْوَانًا خَطَأً: «مَاتَ أَلْفِرِد نُوبِل تَاجِرُ الْمَوْتِ الَّذِي أَصْبَحَ
غَنِيًّا مِنْ خِلَالِ إِيجَادِ طُرُقٍ لِقَتْلِ الْمَزِيدِ مِنَ النَّاسِ».

شَعَرَ نُوبِلُ بِالذَّنْبِ وَ بِخَيْبَةِ الْأَمَلِ مِنْ هٰذَا الْعُنْوَانِ، وَ
بَقِيَ حَزِينًا وَ خَافَ أَنْ يَذْكُرَهُ النَّاسُ بِالسُّوءِ بَعْدَ مَوْتِهِ.
لِذٰلِكَ فَقَدْ بَنَى مُؤَسَّسَةً لِمَنْحِ الْجَوَائِزِ الشَّهِيرَةِ بِاسْمِ
«جَائِزَةِ نُوبِل». وَ مَنَحَ ثَرْوَتَهُ لِشِرَاءِ الْجَوَائِزِ الذَّهَبِيَّةِ
لِكَيْ يُصَحِّحَ خَطَأَهُ.

تُمْنَحُ هٰذِهِ الْجَائِزَةُ فِي كُلِّ سَنَةٍ إِلَى مَنْ يُفِيدُ الْبَشَرِيَّةَ
فِي مَجَالَاتٍ حَدَّدَهَا، وَ هِيَ مَجَالَاتُ السَّلَامِ وَ الْكِيمِيَاءِ وَ
الْفِيزِيَاءِ وَ الطِّبِّ وَ الْأَدَبِ وَ ...

وَلٰكِنْ هَلْ تُعْطَى الْجَوَائِزُ الْيَوْمَ لِمَنْ هُوَ أَهْلٌ لِذٰلِكَ؟!

لِكُلِّ اخْتِرَاعٍ عِلْمِيٍّ وَ ابْتِكَارٍ فِي التِّقْنِيَّةِ وَجْهَانِ:
{\Ar\bfseries وَجْهٌ نَافِعٌ، وَ وَجْهٌ مُضِرٌّ.}
وَ عَلَى الْإِنْسَانِ الْعَاقِلِ أَنْ يَسْتَفِيدَ مِنَ الْوَجْهِ النَّافِعِ.
\end{JNass}

% ============================================================================
%  ۲) ترجمهٔ روان
% ============================================================================
\Jsec{اَلتَّرْجَمَةُ \,\textbar\, ترجمهٔ روانِ درس}

\begin{JTarjome}[ترجمهٔ بندِ اوّل]
در سالِ هزار و هشتصد و سی و چهار در کشورِ «سوئد» کودکی زاده شد که «آلفرد
نوبل» نامیده شد. پدرش کارخانه‌ای برای ساختِ مادهٔ «نیتروگلیسیرین»، مایعِ
سریع‌الانفجار، ساخته بود؛ هرچند با حرارتِ اندک.

آلفرد از کوچکی به این ماده علاقه‌مند شد (اهتمام ورزید) و تلاشگرانه برای
بهینه‌سازیِ آن کار کرد تا از انفجارِ آن جلوگیری کند. آزمایشگاهِ کوچکی برای
انجامِ آزمایش‌هایش ساخت، ولی متأسفانه آزمایشگاه منفجر شد و بر سرِ برادرِ
کوچک‌ترش فروریخت و او را کشت. این رویداد، ارادهٔ او را ضعیف نکرد؛ پس کارش را
با پشتکار ادامه داد تا توانست مادهٔ «دینامیت» را اختراع کند؛ ماده‌ای که جز به
ارادهٔ انسان منفجر نمی‌شود. پس از آن‌که دینامیت را اختراع کرد، رؤسای شرکت‌های
ساختمانی و معادن و نیروهای مسلح — که برای استفاده از آن مشتاق بودند — به
خریدِ آن روی آوردند؛ پس دینامیت در سراسر جهان گسترش یافت.
\end{JTarjome}

\begin{JTarjome}[ترجمهٔ بندِ دوم]
آلفرد به ساختنِ ده‌ها کارخانه و کارگاه در بیست کشور اقدام کرد و از آن ثروتی
بسیار بزرگ به دست آورد؛ چنان‌که از ثروتمندترینِ ثروتمندانِ جهان شد.

انسان از این ماده بهره برد و کارهای دشوارش را در کندنِ تونل‌ها و شکافتنِ
کانال‌ها و ساختنِ راه‌ها و کندنِ معادن و دگرگون‌کردنِ کوه‌ها و تپه‌ها به
دشت‌هایِ قابلِ کشت آسان کرد.

و از کارهای بزرگِ که به‌واسطهٔ این ماده انجام شد، انفجارِ زمین در کانالِ
«پاناما» به اندازه‌ای از دینامیت بود که به چهل تُن می‌رسید.

و پس از آن‌که نوبل دینامیت را اختراع کرد، جنگ‌ها افزایش یافت و ابزارهای
کشتن و ویرانی با این ماده فراوان شد؛ گرچه هدفِ او از اختراعِ آن، یاریِ انسان
در زمینهٔ آبادانی و ساختمان بود.
\end{JTarjome}

\begin{JTarjome}[ترجمهٔ بندِ سوم]
یکی از روزنامه‌های فرانسوی هنگامِ مرگِ برادرِ دیگرش عنوانی اشتباه چاپ کرد:
«آلفرد نوبل، تاجرِ مرگ، مُرد؛ کسی که از راهِ یافتنِ شیوه‌هایی برای کشتنِ
انسان‌های بیشتر ثروتمند شده بود».

نوبل از این عنوان احساسِ گناه و ناامیدی کرد و غمگین ماند و ترسید که مردم پس
از مرگش او را به بدی یاد کنند. پس بنیادی برای اعطای جایزه‌های معروف به نام
«جایزهٔ نوبل» تأسیس کرد و ثروتِ خود را برای خریدِ جایزه‌های طلایی وقف کرد تا
خطای خود را اصلاح کند.

این جایزه هر سال به کسی داده می‌شود که به بشریت در زمینه‌هایی که مشخص کرده
سود رساند؛ و آن زمینه‌ها عبارت‌اند از: صلح، شیمی، فیزیک، پزشکی، ادبیات و ...

ولی آیا امروز این جوایز به شایستگان اهدا می‌شود؟!

برای هر اختراعِ علمی و نوآوریِ فناورانه دو چهره است:
{\bfseries چهره‌ای سودمند، و چهره‌ای زیان‌بار.}
و بر انسانِ خردمند است که از چهرهٔ سودمند بهره گیرد.
\end{JTarjome}

% ============================================================================
%  ۳) واژه‌نامهٔ درس
% ============================================================================
\Jsec{اَلْمُعْجَمُ \,\textbar\, واژه‌نامهٔ درس}

\begin{JMTable}[اَلْمُعْجَمُ \,\textbar\, واژه‌نامهٔ درسِ دوم]
\JW{أَجْرَى}{اجرا کرد}{مضارع: يُجْرِي ← لِيُجْرِيَ: تا اجرا کند}
\JW{أَضْعَفَ}{ضعیف کرد}{}
\JW{اَلْإِعْمَار}{آباد کردن}{أَعْمَرَ ، يُعْمِرُ}
\JW{أَفَادَ}{سود رساند}{مضارع: يُفِيدُ}
\JW{أَقْبَلَ عَلَى}{به... روی آورد}{}
\JW{اَلْأَنْحَاء}{سمت‌ها، سوها}{مفرد: اَلنَّحْو}
\JW{اَلْأَنْفَاق}{تونل‌ها}{مفرد: اَلنَّفَق}
\JW{اِنْهَدَمَ}{ویران شد}{}
\JW{اِهْتَمَّ}{اهتمام ورزید}{مضارع: يَهْتَمُّ}
\JW{اَلْأَهْل}{شایسته}{}
\JW{اَلتَّحْوِيل}{دگرگونی}{}
\JW{اَلتَّطْوِير}{بهینه‌سازی}{}
\JW{اَلتِّقْنِيَّة}{فناوری (تکنیک)}{}
\JW{اَلتِّلَال}{تپه‌ها}{مفرد: اَلتَّلّ}
\JW{تَمَّ}{انجام شد، کامل شد}{مضارع: يَتِمُّ}
\JW{جُلْبُ شَعِيرَةٍ}{پوستِ جو}{}
\JW{حَدَّدَ}{مشخص کرد}{}
\JW{خَيْبَةُ الْأَمَلِ}{ناامیدی}{\textbf{≠} اَلرَّجَاء}
\JW{اَلدَّؤُوب}{با پشتکار}{}
\JW{اَلسُّوَيْد}{سوئد}{}
\JW{سَهَّلَ}{آسان کرد}{\textbf{≠} صَعَّبَ}
\JW{اَلسُّهُول}{دشت‌ها}{مفرد: اَلسَّهْل}
\JW{اَلشَّعِير}{جو}{}
\JW{اَلشَّقّ}{شکافتن}{شَقَّ ، يَشُقُّ}
\JW{اَلصَّالِحَةُ لِلزِّرَاعَةِ}{قابلِ کشت}{}
\JW{اَلصَّبِيّ}{کودک، پسر}{جمع: اَلصِّبْيَان}
\JW{صَحَّحَ}{تصحیح کرد}{}
\JW{اَلطُّنّ}{تُن}{جمع: اَلْأَطْنَان}
\JW{اَلْفَرَنْسِيَّة}{زبانِ فرانسوی}{}
\JW{اَلْفِيزِيَاء}{فیزیک}{}
\JW{قَنَاةُ بَنَمَا}{کانالِ پاناما}{}
\JW{اَلْقَنَوَات}{کانال‌ها}{مفرد: اَلْقَنَاة}
\JW{كَسَبَ}{به دست آورد}{}
\JW{اَلْمَجَال}{زمینه}{جمع: اَلْمَجَالَات}
\JW{اَلْمَنَاجِم}{معادن}{مفرد: اَلْمَنْجَم}
\JW{نَشَرَ}{پخش کرد}{}
\end{JMTable}

% ============================================================================
%  ۴) قواعدِ درس
% ============================================================================
\Jsec{إِعْلَمُوا \,\textbar\, آموزشِ قواعد}

\Jsubsec{اَلْحَال \,\textbar\, قیدِ حالت}

آیا ترجمهٔ این سه جمله یکسان است؟
\begin{JMesal}
\begin{tabular}{@{}p{0.52\linewidth}@{\hspace{0.6em}}p{0.42\linewidth}@{}}
\Ar (أ) رَأَيْتُ وَلَدًا مَسْرُورًا. & (پسرِ خوشحالی را دیدم.)\\[2pt]
\Ar (ب) رَأَيْتُ الْوَلَدَ الْمَسْرُورَ. & (پسرِ خوشحال را دیدم.)\\[2pt]
\Ar (ج) رَأَيْتُ الْوَلَدَ، مَسْرُورًا. & (پسر را خوشحال \emph{دیدم}.)\\
\end{tabular}
\end{JMesal}

\emph{مَسْرُورًا} در جملهٔ (أ) و \emph{الْمَسْرُورَ} در جملهٔ (ب) چه نقشی
دارند؟ آیا \emph{مَسْرُورًا} در جملهٔ (ج) نیز همان نقش را دارد؟

در جملهٔ اوّل، «مَسْرُورًا» \textbf{صفتِ} «وَلَدًا» و در جملهٔ دوم
«الْمَسْرُورَ» صفتِ «الْوَلَدَ» است؛ ولی در جملهٔ سوم «مَسْرُورًا»
\textbf{حالتِ} «الْوَلَدَ» هنگامِ وقوعِ فعل را نشان می‌دهد:

\begin{JMesal}
\noindent
\begin{tabular}{@{}p{0.48\linewidth}@{\hspace{0.02\linewidth}}p{0.48\linewidth}@{}}
\Ar ذَهَبَتِ الْبِنْتُ الْفَرِحَةُ. &
\Ar ذَهَبَتِ الْبِنْتُ، فَرِحَةً.\\[2pt]
{\small الْبِنْتُ: موصوف} & {\small فَرِحَةً: حال (قیدِ حالت)}\\
{\small الْفَرِحَةُ: صفت} & {\small}\\
\end{tabular}
\end{JMesal}

\begin{JQaede}[JCompact]{قاعدهٔ اَلْحَال (قیدِ حالت)}
برخی واژه‌ها در جمله، \textbf{حالتِ یک اسم را هنگامِ وقوعِ فعل} نشان می‌دهند؛
به چنین واژه‌هایی در فارسی «\textbf{قیدِ حالت}» و در عربی «\textbf{حَال}»
می‌گویند.

\begin{itemize}
  \item حال، \textbf{منصوب} است و بر صاحبِ حال (مرجعِ حال) بازمی‌گردد.
  \item در ترجمه، اغلب با «\textbf{در حالی‌که}»، «\textbf{به‌صورتِ}»،
        «\textbf{با حالتِ}» یا صفتِ حالی می‌آید: «خوشحال \emph{دیدم}»،
        «با کوششِ \emph{فراوان}».
  \item در زبانِ عربی بسیاری از اوقات، قیدِ حالت در \textbf{انتهایِ جمله}
        می‌آید.
\end{itemize}

مثال‌ها:
\begin{itemize}
  \item اِشْتَغَلَ مَنْصُورٌ فِي الْمَزْرَعَةِ نَشِيطًا. (نَشِيطًا: حال) —
        منصور در مزرعه با کوششِ فراوان کار کرد.
  \item وَقَفَ الْمُهَنْدِسُ الشَّابُّ فِي الْمَصْنَعِ مُبْتَسِمًا. — مهندسِ
        جوان در کارخانه در حالی‌که لبخند می‌زد ایستاد.
  \item رَأَيْتُ الطِّفْلَتَيْنِ مُبْتَسِمَتَيْنِ. — دو کودک را خندان دیدم.
  \item اَللَّاعِبُونَ الْإِيرَانِيُّونَ رَجَعُوا مِنَ الْمُسَابَقَةِ
        مُبْتَسِمِينَ. — بازیکنانِ ایرانی از مسابقه خندان بازگشتند.
\end{itemize}
\end{JQaede}

\begin{JDam}[دامِ آزمون — تفاوتِ حال و صفت]
هر دو «مَسْرُورًا» و «الْمَسْرُورَ» می‌توانند منصوب باشند؛ آزمون با \textbf{تعریف
و تنکری} آن‌ها را جدا می‌کند:
\begin{itemize}
  \item \textbf{صفت:} بعد از موصوفِ \emph{معرفه} با «ال» می‌آید و مطابقِ اوست:
        رَأَيْتُ الْوَلَدَ \textbf{الْمَسْرُورَ}. (پسرِ خوشحال را دیدم — صفتِ
        همیشگی و ثابت)
  \item \textbf{حال:} نکره است و حالتِ لحظه‌ای را می‌رساند:
        رَأَيْتُ الْوَلَدَ، \textbf{مَسْرُورًا}. (پسر را خوشحال \emph{دیدم} —
        در همان لحظه)
\end{itemize}
(تشخیصِ مرجعِ قیدِ حالت و مطابقتِ آن با مرجعش از نظرِ عدد و جنس، از اهدافِ
آموزشیِ کتاب درسی نیست؛ ولی شناختِ «حال» برای ترجمه لازم است.)
\end{JDam}

\Jsubsec{جملهٔ حالیّه با واوِ حالیّه}

گاهی قیدِ حالت به‌صورتِ \textbf{جملهٔ اسمیّه} همراه با حرفِ «\textbf{واوِ
حالیّه}» و به دنبالِ آن یک ضمیر می‌آید؛ مثال:

\begin{JMesal}
\Ar (أ) رَأَيْتُ الْفَلَّاحَ وَ هُوَ يَجْمَعُ الْمَحْصُولَ.

کشاورز را دیدم، در حالی‌که محصول را جمع می‌کرد.

\Ar (ب) أُشَاهِدُ قَاسِمًا وَ هُوَ جَالِسٌ بَيْنَ الشَّجَرَتَيْنِ.

قاسم را می‌بینم، در حالی‌که میانِ دو درخت نشسته است.
\end{JMesal}

\begin{JNokte}[نکتهٔ ترجمه]
در جملهٔ (أ)، قیدِ حالتِ جمله‌ای به‌صورتِ \textbf{ماضیِ استمراری} ترجمه
می‌شود: «در حالی‌که \emph{جمع می‌کرد}». پس حضورِ واوِ حالیّه + ضمیر، به‌تنهایی
معیارِ مضارع بودنِ ترجمه نیست؛ بافتِ جمله را ببین!
\end{JNokte}

\begin{JGood}[تحلیلِ صرفی و اعراب]
به بیانِ ویژگی‌های دستوری (ساختاری)ِ واژه‌های بیرون از جمله، در زبانِ عربی
«\textbf{التَّحْلِيلُ الصَّرْفِيُّ}» (تحلیلِ صرفی) می‌گویند؛ و به ذکرِ نقشِ
دستوریِ واژه — یعنی جایگاهِ آن در جمله — و علامتِ انتهایِ واژه،
«\textbf{إِعْرَاب}» می‌گویند. (در تمرینِ پنجم این درس با هر دو کار داریم.)
\end{JGood}

% ============================================================================
%  ۵) خودآزمایی (بدونِ پاسخ)
% ============================================================================
\Jsec{اِخْتَبِرْ نَفْسَكَ \,\text| خودآزماییِ درس}

\begin{JTest}[خودآزماییِ ١ \,\textbar\, تعیینِ حال]
عَيِّنِ «الْحَالَ» فِي الْجُمَلِ التَّالِيَةِ (حال را در جملاتِ زیر مشخص کن):

\begin{enumerate}
  \item وَصَلَ الْمُسَافِرَانِ إِلَى الْمَطَارِ مُتَأَخِّرَيْنِ وَ رَكِبَا
        الطَّائِرَةَ.
  \item تَجْتَهِدُ الطَّالِبَةُ فِي أَدَاءِ وَاجِبَاتِهَا رَاضِيَةً وَ
        تُسَاعِدُ أُمَّهَا.
  \item يُشَجِّعُ الْمُتَفَرِّجُونَ فَرِيقَهُمُ الْفَائِزَ فَرِحِينَ.
  \item اَلطَّالِبَتَانِ تَقْرَآنِ دُرُوسَهُمَا مُجِدَّتَيْنِ.
\end{enumerate}
\end{JTest}

\begin{JTest}[خودآزماییِ ٢ \,\text| ترجمهٔ آیات و تعیینِ حال]
تَرْجِمِ الْآيَاتِ الْكَرِيمَةَ، ثُمَّ عَيِّنِ «الْحَالَ»:

\begin{enumerate}
  \item ﴿... وَ خُلِقَ الْإِنْسَانُ ضَعِيفًا﴾ اَلنِّسَاء: ٢٨
  \item ﴿وَ لا تَهِنُوا\FNmark{١} وَ لا تَحْزَنُوا وَ أَنْتُمُ الْأَعْلَوْنَ
        ...﴾ آلِ عِمْرَان: ١٣٩
        \FNfoot{١}{لا تَهِنُوا: سست نشوید ← وَهَنَ: سست شد}
  \item ﴿كَانَ النَّاسُ أُمَّةً وَاحِدَةً فَبَعَثَ اللّٰهُ النَّبِيِّينَ
        مُبَشِّرِينَ ...﴾ اَلْبَقَرَة: ٢١٣
  \item ﴿يَا أَيَّتُهَا النَّفْسُ الْمُطْمَئِنَّةُ ارْجِعِي إِلَى رَبِّكِ
        رَاضِيَةً مَرْضِيَّةً\FNmark{٢}﴾ اَلْفَجْر: ٢٧ و ٢٨
        \FNfoot{٢}{مَرْضِيَّةً: با خشنودی و خداپسندیدگی}
  \item ﴿إِنَّمَا وَلِيُّكُمُ اللّٰهُ وَ رَسُولُهُ وَ الَّذِينَ آمَنُوا
        الَّذِينَ يُقِيمُونَ الصَّلَاةَ وَ يُؤْتُونَ\FNmark{٣} الزَّكَاةَ وَ
        هُمْ رَاكِعُونَ﴾ اَلْمَائِدَة: ٥٥
        \FNfoot{٣}{يُؤْتُونَ: می‌دهند ← آتَى ، يُؤْتِي}
\end{enumerate}
\end{JTest}

% ============================================================================
%  ۶) تمرین‌های کتاب (بدونِ پاسخ)
% ============================================================================
\Jsec{اَلتَّمَارِينُ \,\text| تمرین‌های کتاب}

\begin{JTamrin}{١}{جست‌وجو در معجم}
اِبْحَثْ فِي مُعْجَمِ الدَّرْسِ عَنْ كَلِمَةٍ مُنَاسِبَةٍ لِلتَّوْضِيحَاتِ
التَّالِيَةِ (واژهٔ مناسبِ هر تعریف را از معجمِ درس بیاور):

\begin{enumerate}
  \item أُسْلُوبٌ أَوْ فَنٌّ فِي إِنْجَازِ\FNmark{١} عَمَلٍ، أَوْ طَرِيقٌ
        يَخْتَصُّ بِمِهْنَةٍ، أَوْ عِلْمُ الصِّنَاعَةِ الْحَدِيثَةِ.
        \JdotsW{3.5cm} \FNfoot{١}{إِنْجَاز: انجام}
  \item مَمَرٌّ تَحْتَ الْأَرْضِ أَوْ فِي الْجَبَلِ، طُولُهُ أَكْثَرُ مِنْ
        عَرْضِهِ، لَهُ مَدْخَلٌ وَ مَخْرَجٌ. \JdotsW{3.5cm}
  \item مِنْطَقَةٌ مُرْتَفِعَةٌ فَوْقَ سَطْحِ الْأَرْضِ، أَصْغَرُ مِنَ
        الْجَبَلِ. \JdotsW{3.5cm}
  \item مَكَانُ الذَّهَبِ وَ الْفِضَّةِ وَ النُّحَاسِ وَ نَحْوِهَا فِي
        الْأَرْضِ. \JdotsW{3.5cm}
  \item اَلَّذِي يَسْعَى فِي إِنْجَازِ عَمَلِهِ، وَ لا يَشْعُرُ بِالتَّعَبِ.
        \JdotsW{3.5cm}
  \item وَزْنٌ يُعَادِلُ أَلْفَ كِيلُوغْرَام. \JdotsW{3.5cm}
\end{enumerate}
\end{JTamrin}

\begin{JTamrin}{٢}{خطبهٔ ٢٢٤ نهج‌البلاغه}
اِقْرَأِ النَّصَّ التَّالِيَ؛ ثُمَّ عَيِّنْ تَرْجَمَةَ الْكَلِمَاتِ
الْحَمْرَاءِ، وَ عَيِّنِ الْمَطْلُوبَ مِنْكَ.

\begin{JNass}[اَلْخُطْبَةُ الْمِئَتَانِ وَ الرَّابِعَةُ وَ الْعِشْرُونَ
مِنْ نَهْجِ الْبَلَاغَة]
\Ar ... وَ اللّٰهِ، لَوْ أُعْطِيتُ \jr{الْأَقَالِيمَ السَّبْعَةَ} بِمَا تَحْتَ
\jr{أَفْلَاكِهَا} عَلَى أَنْ أَعْصِيَ اللّٰهَ فِي نَمْلَةٍ \jr{أَسْلُبُهَا}
\jr{جُلْبَ} \jr{شَعِيرَةٍ}، مَا فَعَلْتُ، وَ إِنَّ دُنْيَاكُمْ عِنْدِي
\jr{لَأَهْوَنُ} مِنْ وَرَقَةٍ فِي فَمِ \jr{جَرَادَةٍ} \jr{تَقْضَمُهَا}. مَا
لِعَلِيٍّ وَ لِنَعِيمٍ يَفْنَى وَ لَذَّةٍ لا تَبْقَى؟! نَعُوذُ بِاللّٰهِ مِنْ
\jr{سُبَاتِ الْعَقْلِ} وَ \jr{قُبْحِ الزَّلَلِ} وَ بِهِ \jr{نَسْتَعِينُ}.

(كَلِمَاتُ الْحَمْرَاءِ: لَوْ أُعْطِيتُ ، أَسْلُبُ ، جُلْبَ ، شَعِيرَةٍ ،
أَهْوَنُ ، جَرَادَةٍ ، تَقْضَمُ ، نَعِيمٍ ، سُبَاتِ الْعَقْلِ ، قُبْحِ
الزَّلَلِ ، نَسْتَعِينُ)
\end{JNass}

\begin{JTarjome}[خطبهٔ دویست و بیست و چهارمِ نهج‌البلاغه]
به خدا، اگر سرزمین‌های هفتگانه با هرچه زیرِ آسمان‌هایشان هست به من داده می‌شد
تا دربارهٔ مورچه‌ای از خدا نافرمانی کنم و پوستِ جویی را از آن به زور بگیرم،
انجام نمی‌دادم؛ و بی‌گمان دنیای شما نزدِ من از برگِی در دهانِ ملخِی که آن را
می‌جَوَد، پست‌تر است. علی را چه کار با نعمتی که نابود می‌شود و لذّتی که
نمی‌ماند؟! از به‌خواب‌رفتنِ خرد و زشتیِ لغزش به خدا پناه می‌بریم و از او یاری
می‌جوییم.
\end{JTarjome}

پاورقی‌ها: ١- سَلَبَ: أَخَذَ مِنْهُ قَهْرًا (به زور گرفت) ؛ ٢- اَلْجُلْب:
قِشْرُ النَّبَاتِ وَ خَشَبُهُ (پوستِ گیاه و چوبِ آن) ؛ ٣- اَلشَّعِيرَة:
حَبَّةُ نَبَاتٍ (دانهٔ گیاهی) ؛ ٤- اَلْأَهْوَن: اَلْأَحْقَر (پست‌تر) ؛
٥- اَلْجَرَاد: حَشَرَةٌ مُضِرَّةٌ تَأْكُلُ الْمَحَاصِيلَ الزِّرَاعِيَّةَ وَ
النَّبَاتَاتِ ؛ ٦- قَضَمَ: كَسَرَ بِأَطْرَافِ أَسْنَانِهِ (جوید) ؛
٧- اَلنَّعِيم: اَلْمَال ، اَلْجَنَّة ؛ ٨- اَلسُّبَات: اَلنَّوْمُ
الْخَفِيفُ (چرتِ سبک) ؛ ٩- اِسْتَعَانَ: طَلَبَ الْمُسَاعَدَةَ.

پرسش‌ها:
\begin{enumerate}
  \item إعْرَابُ هٰذِهِ الْكَلِمَاتِ:\quad نَمْلَةٍ \JdotsW{2.6cm} جُلْبَ
        \JdotsW{2.6cm} اَلْعَقْلِ \JdotsW{2.6cm} جَرَادَةٍ \JdotsW{2.6cm}
  \item نَوْعُ فِعْلِ «فَعَلْتُ» وَ صِيغَتُهُ: \JdotsW{6cm}
  \item نَوْعُ فِعْلِ «تَقْضَمُ» وَ صِيغَتُهُ: \JdotsW{6cm}
  \item جَمْعَيْنِ لِلتَّكْسِيرِ فِي النَّصِّ: \JdotsW{6cm}
  \item اَلْفِعْلَ الْمَجْهُولَ فِي النَّصِّ: \JdotsW{6cm}
  \item عَدَدَ الْأَفْعَالِ فِي النَّصِّ: \JdotsW{6cm}
\end{enumerate}
\end{JTamrin}

\begin{JTamrin}{٣}{واژهٔ زائد}
ضَعْ فِي الدَّائِرَةِ الْعَدَدَ الْمُنَاسِبَ. «كَلِمَةٌ وَاحِدَةٌ زَائِدَةٌ»
(در هر سطر، شمارهٔ واژه‌ای که با توضیحِ آن سطر هماهنگ نیست را دایره بزن):

\begin{enumerate}
  \item اَلْقَنَاةُ (١)\quad تُرَابٌ مُخْتَلِطٌ بِالْمَاءِ، وَ قَدْ يُسَمَّى
        بِذٰلِكَ وَ إِنْ زَالَتْ\FNmark{١} عَنْهُ الرُّطُوبَةُ.
        \FNfoot{١}{زَالَ: سپری و نابود شد}
  \item اَلْفِيزِيَاءُ (٢)\quad شَرِيطٌ\FNmark{٢} يَسْتَعْمِلُهُ رُكَّابُ
        الطَّائِرَةِ وَ السَّيَّارَاتِ لِلنَّجَاةِ مِنَ الْخَطَرِ.
        \FNfoot{٢}{شَرِيط: نوار}
  \item اَلطِّينُ (٣)\quad عِلْمٌ يَبْحَثُ عَنْ خَصَائِصِ الْمَوَادِّ وَ
        الظَّوَاهِرِ الطَّبِيعِيَّةِ وَ الطَّاقَةِ.
  \item اَلتِّلَالُ (٤)\quad حَشَرَةٌ تَأْكُلُ الْمَحَاصِيلَ الزِّرَاعِيَّةَ
        تَسْتَطِيعُ أَنْ تَقْفِزَ مِتْرًا وَاحِدًا.
  \item حِزَامُ الْأَمَانِ (٥)\quad نَهْرٌ وَاسِعٌ أَوْ ضَيِّقٌ\FNmark{٣}
        لِحَرَكَةِ الْمِيَاهِ مِنْ مَكَانٍ إِلَى آخَرَ.
        \FNfoot{٣}{اَلضَّيِّق: تنگ}
  \item اَلْجَرَادَةُ (٦)
\end{enumerate}
\end{JTamrin}

\begin{JTamrin}{٤}{ترجمهٔ دانستنی‌ها}
لِلتَّرْجَمَةِ. (هَلْ تَعْلَمُ أَنَّ ... ؟ — «آیا می‌دانی که...؟» را کامل
ترجمه کن):

\begin{enumerate}
  \item ... تَلَفُّظَ «گ» وَ «چ» وَ «ژ» مَوْجُودٌ فِي اللَّهَجَاتِ
        الْعَرَبِيَّةِ الدَّارِجَةِ\FNmark{١} كَثِيرًا؟ \Jdots
        \FNfoot{١}{اَلدَّارِجَة: عامیانه، رایج}
  \item ... الْمَغُولَ اسْتَطَاعُوا أَنْ يَهْجُمُوا عَلَى الصِّينِ عَلَى
        رَغْمِ بِنَاءِ سُورٍ\FNmark{٢} عَظِيمٍ حَوْلَهَا؟ \Jdots
        \FNfoot{٢}{اَلسُّور: دیوار}
  \item ... الْحُوتَ يُصَادُ لِاسْتِخْرَاجِ الزَّيْتِ مِنْ كَبِدِهِ لِصِنَاعَةِ
        مَوَادِّ التَّجْمِيلِ\FNmark{٣}؟ \Jdots
        \FNfoot{٣}{مَوَادُّ التَّجْمِيلِ: موادّ آرایشی}
  \item ... الْخُفَّاشَ هُوَ الْحَيَوَانُ اللَّبُونُ الْوَحِيدُ الَّذِي
        يَقْدِرُ عَلَى الطَّيَرَانِ\FNmark{٤}؟ \Jdots
        \FNfoot{٤}{اَلطَّيَرَان: پرواز، پرواز کردن}
  \item ... عَدَدَ النَّمْلِ فِي الْعَالَمِ يَفُوقُ عَدَدَ الْبَشَرِ بِمِلْيُونِ
        مَرَّةٍ تَقْرِيبًا؟ \Jdots
  \item ... طَيْسَفُون قُرْبَ بَغْدَادَ كَانَتْ عَاصِمَةَ\FNmark{٥}
        السَّاسَانِيِّينَ؟ \Jdots
        \FNfoot{٥}{اَلْعَاصِمَة: پایتخت — جمع: اَلْعَوَاصِم}
  \item ... حَجْمَ دُبِّ الْبَانْدَا\FNmark{٦} عِنْدَ الْوِلَادَةِ أَصْغَرُ
        مِنَ الْفَأْرِ؟ \Jdots \FNfoot{٦}{دُبُّ الْبَانْدَا: خرسِ پاندا}
  \item ... الزَّرَافَةَ بَكْمَاءُ لَيْسَتْ لَهَا أَحْبَالٌ
        صَوْتِيَّةٌ\FNmark{٧}؟ \Jdots
        \FNfoot{٧}{اَلْأَحْبَالُ الصَّوْتِيَّةُ: تارهای صوتی — اَلْأَحْبَال:
        جمع ، اَلْحَبْل: مفرد}
  \item ... وَرَقَةَ الزَّيْتُونِ رَمْزُ\FNmark{٨} السَّلَامِ؟ \Jdots
        \FNfoot{٨}{اَلرَّمْز: نماد، سمبل — جمع: اَلرُّمُوز}
\end{enumerate}
\end{JTamrin}

\begin{JTamrin}{٥}{تحلیلِ صرفی و اعراب}
عَيِّنِ الصَّحِيحَ فِي التَّحْلِيلِ الصَّرْفِيِّ وَ الْإِعْرَابِ لِمَا أُشِيرَ
إِلَيْهِ بِخَطٍّ (گزینهٔ صحیحِ تحلیل و اعرابِ واژه‌های زیرخط‌دار را برگزین):

\begin{center}
\Ar اَلْعُمَّالُ الْمُجْتَهِدُونَ يَشْتَغِلُونَ فِي الْمَصْنَعِ.\quad
\Ar نَجَحَتِ الطَّالِبَاتُ فِي الاِمْتِحَانِ.
\end{center}

\begin{multicols}{2}
\begin{enumerate}
  \item اَلْعُمَّالُ\\
        (أ) اِسْمٌ ، جَمْعُ تَكْسِيرٍ ، اِسْمُ فَاعِلٍ ، مُعَرَّفٌ بِأَل ،
        مُعْرَبٌ \textbar\ مُبْتَدَأٌ وَ مَرْفُوعٌ بِالضَّمَّةِ\\
        (ب) اِسْمُ مُبَالَغَةٍ ، جَمْعٌ مُكَسَّرٌ وَ مُفْرَدُهُ «اَلْعَامِلُ»
        \textbar\ فَاعِلٌ وَ مَرْفُوعٌ بِالضَّمَّةِ
  \item اَلْمُجْتَهِدُونَ\\
        (أ) اِسْمُ فَاعِلٍ ، جَمْعٌ مُذَكَّرٌ سَالِمٌ ، مُعَرَّفٌ بِأَل
        \textbar\ صِفَةٌ وَ مَرْفُوعَةٌ بِالتَّبَعِيَّةِ لِمَوْصُوفِهَا\\
        (ب) اِسْمُ مَفْعُولٍ ، مُثَنَّى ، مُذَكَّرٌ ، نَكِرَةٌ ، مَبْنِيٌّ
        \textbar\ مُضَافٌ إِلَيْهِ وَ مَجْرُورٌ بِـ«و» فِي «ـونَ»
  \item يَشْتَغِلُونَ\\
        (أ) فِعْلٌ ، جَمْعٌ مُذَكَّرٌ غَائِبٌ ، ثُلَاثِيٌّ مَزِيدٌ مِنْ بَابِ
        اِفْتِعَال ، مَعْلُومٌ ، لَازِمٌ \textbar\ خَبَرٌ\\
        (ب) فِعْلٌ مُضَارِعٌ ، جَمْعٌ مُذَكَّرٌ مُخَاطَبٌ ، مَجْهُولٌ ،
        ثُلَاثِيٌّ مُجَرَّدٌ ، مُتَعَدٍّ \textbar\ فَاعِلٌ
  \item اَلْمَصْنَعِ\\
        (أ) اِسْمُ مَفْعُولٍ ، مُفْرَدٌ ، مُعَرَّفٌ بِأَل \textbar\ مُضَافٌ
        إِلَيْهِ وَ مَجْرُورٌ\\
        (ب) اِسْمُ مَكَانٍ ، مُفْرَدٌ ، مُذَكَّرٌ ، مَعْرِفَةٌ \textbar\
        مَجْرُورٌ بِحَرْفِ جَرٍّ (فِي الْمَصْنَعِ: جَارٌّ وَ مَجْرُورٌ)
  \item نَجَحَتْ\\
        (أ) فِعْلٌ مَاضٍ ، مُفْرَدٌ مُذَكَّرٌ مُخَاطَبٌ ، ثُلَاثِيٌّ مَزِيدٌ
        مِنْ بَابِ إِفْعَال ، مَجْهُولٌ\\
        (ب) فِعْلٌ مَاضٍ ، مُفْرَدٌ مُؤَنَّثٌ غَائِبٌ ، ثُلَاثِيٌّ مُجَرَّدٌ ،
        مَعْلُومٌ ، لَازِمٌ
  \item اَلطَّالِبَاتُ\\
        (أ) مَصْدَرٌ ، جَمْعٌ مُذَكَّرٌ سَالِمٌ ، مُعَرَّفٌ بِأَل \textbar\
        مُبْتَدَأٌ وَ مَرْفُوعٌ بِالضَّمَّةِ\\
        (ب) اِسْمُ فَاعِلٍ ، جَمْعٌ مُؤَنَّثٌ سَالِمٌ ، مَعْرِفَةٌ \textbar\
        فَاعِلٌ وَ مَرْفُوعٌ بِالضَّمَّةِ
  \item اَلاِمْتِحَانِ\\
        (أ) مَصْدَرٌ مِنْ بَابِ اِفْتِعَالٍ ، مُفْرَدٌ ، مُذَكَّرٌ ، مُعَرَّفٌ
        بِأَل \textbar\ مَجْرُورٌ بِحَرْفِ جَرٍّ (فِي الاِمْتِحَانِ: جَارٌّ
        وَ مَجْرُورٌ)\\
        (ب) اِسْمُ تَفْضِيلٍ ، مُفْرَدٌ ، مُذَكَّرٌ ، مَعْرِفَةٌ بِالْعَلَمِيَّةِ
        \textbar\ صِفَةٌ وَ مَجْرُورَةٌ بِالتَّبَعِيَّةِ لِمَوْصُوفِهَا
\end{enumerate}
\end{multicols}
\end{JTamrin}

\begin{JTamrin}{٦}{ترجمهٔ خانوادهٔ واژه‌ها}
لِلتَّرْجَمَةِ (ترجمهٔ هر جمله را با توجّه به معنای خانوادهٔ واژه بنویس):

\begin{multicols}{2}
\begin{enumerate}
  \item \textbf{كَتَبَ: نوشت}\\
        كَتَبَ التَّمْرِينَ: \JdotsW{3.4cm}\\
        لِمَ لا تَكْتُبِينَ دَرْسَكِ؟ \JdotsW{3.4cm}\\
        لَمْ تَكْتُبِي شَيْئًا: \JdotsW{3.4cm}
  \item \textbf{تَكَاتَبَ: نامه‌نگاری کرد}\\
        اَلصَّدِيقَانِ تَكَاتَبَا: \JdotsW{3.4cm}\\
        أَيُّهَا الصَّدِيقَانِ، تَكَاتَبَا: \JdotsW{3.4cm}\\
        تَكَاتَبَ الزَّمِيلَانِ: \JdotsW{3.4cm}
  \item \textbf{مَنَعَ: بازداشت، منع کرد}\\
        مُنِعْتُ عَنِ الْمَوَادِّ السُّكَّرِيَّةِ: \JdotsW{3.4cm}\\
        لا تَمْنَعْنَا عَنِ الْخُرُوجِ: \JdotsW{3.4cm}\\
        شَاهَدْنَا مَانِعًا بِالطَّرِيقِ: \JdotsW{3.4cm}
  \item \textbf{اِمْتَنَعَ: خودداری کرد}\\
        كَانَ الْحَارِسُ قَدِ امْتَنَعَ عَنِ النَّوْمِ: \JdotsW{3.4cm}\\
        لا تَمْتَنِعُوا عَنِ الْمُطَالَعَةِ: \JdotsW{3.4cm}\\
        لَنْ نَمْتَنِعَ عَنِ الْخُرُوجِ: \JdotsW{3.4cm}
  \item \textbf{عَمِلَ: کار کرد، عمل کرد}\\
        لِمَ مَا عَمِلْتُمْ بِوَاجِبَاتِكُمْ؟ \JdotsW{3.4cm}\\
        أَ تَعْمَلُونَ فِي الْمَصْنَعِ؟ \JdotsW{3.4cm}\\
        اَلْعُمَّالُ مَشْغُولُونَ بِالْعَمَلِ: \JdotsW{3.4cm}
  \item \textbf{عَامَلَ: رفتار کرد}\\
        إلٰهِي، عَامِلْنَا بِفَضْلِكَ: \JdotsW{3.4cm}\\
        إلٰهِي، لا تُعَامِلْنَا بِعَدْلِكَ: \JdotsW{3.4cm}\\
        كَانُوا يُعَامِلُونَنَا جَيِّدًا: \JdotsW{3.4cm}
  \item \textbf{ذَكَرَ: یاد کرد}\\
        قَدْ ذَكَرَ الْمُؤْمِنُ رَبَّهُ: \JdotsW{3.4cm}\\
        ذُكِرْتَ بِالْخَيْرِ: \JdotsW{3.4cm}\\
        يَذْكُرُ الْأُسْتَاذُ تَلَامِيذَهُ الْقُدَمَاءَ: \JdotsW{3.4cm}
  \item \textbf{تَذَكَّرَ: به یاد آورد}\\
        جَدِّي وَ جَدَّتِي تَذَكَّرَانِي: \JdotsW{3.4cm}\\
        سَيَتَذَكَّرُنَا الْمُدَرِّسُ: \JdotsW{3.4cm}\\
        لا أَتَذَكَّرُكَ يَا زَمِيلِي: \JdotsW{3.4cm}
\end{enumerate}
\end{multicols}
\end{JTamrin}

\begin{JTamrin}{٧}{تعیینِ حال}
عَيِّنِ «الْحَالَ» فِي الْعِبَارَاتِ التَّالِيَةِ:

\begin{enumerate}
  \item مَنْ عَاشَ بِوَجْهَيْنِ، مَاتَ خَاسِرًا. \JdotsW{5.5cm}
  \item أَقْوَى النَّاسِ مَنْ عَفَا\FNmark{١} عَدُوَّهُ مُقْتَدِرًا.
        \JdotsW{5.5cm} \FNfoot{١}{عَفَا: عفو کرد}
  \item عِنْدَ وُقُوعِ الْمَصَائِبِ تَذْهَبُ الْعَدَاوَةُ سَرِيعَةً.
        \JdotsW{5.5cm}
  \item مَنْ أَذْنَبَ\FNmark{٢} وَ هُوَ يَضْحَكُ، دَخَلَ النَّارَ وَ هُوَ
        يَبْكِي. \JdotsW{5.5cm} \FNfoot{٢}{أَذْنَبَ: گناه کرد}
  \item يَبْقَى الْمُحْسِنُ حَيًّا وَ إِنْ نُقِلَ إِلَى مَنَازِلِ الْأَمْوَاتِ.
        \JdotsW{5.5cm}
  \item إِذَا طَلَبْتَ أَنْ تَنْجَحَ فِي عَمَلِكَ فَقُمْ بِهِ وَحِيدًا وَ لا
        تَتَوَكَّلْ عَلَى النَّاسِ. \JdotsW{5.5cm}
\end{enumerate}
\end{JTamrin}

\begin{JTamrin}{٨}{مفرد و جمع}
عَيِّنِ الصَّحِيحَ فِي الْمُفْرَدِ وَ جَمْعِهِ (درست یا نادرست بودنِ هر جفت را
مشخص کن):

\begin{multicols}{2}
\begin{enumerate}
  \item مَصْنَع ، صِنَاعَات
  \item طَعَام ، مَطَاعِم
  \item ظَاهِرَة ، مَظَاهِر
  \item صَبِيّ ، صِبْيَان
  \item عَظْم ، أَعَاظِم
  \item مَمْلَكَة ، مُلُوك
  \item رَئِيس ، رُؤُوس
  \item قَنَاة ، قَنَوَات
  \item عَامِل ، عُمَّال
  \item حَبْل ، أَحْبَال
  \item طَرِيق ، طُرُق
  \item مَنْجَم ، أَنْجُم
  \item جُرْم ، جَرَائِم
  \item أَدَاة ، أَدَوَات
  \item سَهْل ، سُهُول
  \item عَمَل ، أَعْلَام
  \item نَفَق ، أَنْفَاق
  \item لَحْم ، لُحُوم
  \item نَحْو ، أَنْحَاء
  \item تَلّ ، تِلَال
\end{enumerate}
\end{multicols}
\end{JTamrin}

% ============================================================================
%  ۷) پاسخ‌نامهٔ خودآزمایی و تمرین‌ها
% ============================================================================
\Jsec{اَلْإِجَابَاتُ \,\text| پاسخ‌نامهٔ خودآزمایی و تمرین‌ها}

\begin{JKey}[پاسخ‌نامهٔ خودآزمایی]
\JAnswer{١}{حال‌ها: ۱. مُتَأَخِّرَيْنِ (مسافران با تأخیر به فرودگاه رسیدند) ؛
۲. رَاضِيَةً (دانش‌آموز، تکلیف‌هایش را با خشنودی انجام می‌دهد) ؛
۳. فَرِحِينَ (تماشاگران تیمِ برندهٔ خود را خوشحال تشویق می‌کنند) ؛
۴. مُجِدَّتَيْنِ (دو دانش‌آموز درس‌هایشان را با تلاش می‌خوانند).}

\JAnswer{٢}{%
  ۱. «و انسان ضعیف آفریده شد» — حال: ضَعِيفًا.\\
  ۲. «و سست نشوید و اندوهگین نشوید در حالی‌که شما برترید» — حالِ جمله‌ای:
  وَ أَنْتُمُ الْأَعْلَوْنَ (با واوِ حالیّه).\\
  ۳. «مردم امّتی واحد بودند، پس خدا پیامبران را بشارت‌دهنده برانگیخت» — حال:
  مُبَشِّرِينَ.\\
  ۴. «ای نفسِ مطمئنه! به سوی پروردگارت بازگرد؛ در حالی‌که تو از او خشنودی و
  او از تو خشنود است» — حال: رَاضِيَةً مَرْضِيَّةً.\\
  ۵. «ولیّ شما تنها خدا و پیامبرِ او و کسانی‌اند که ایمان آورده‌اند؛ همانان
  که نماز را برپا می‌دارند و زکات می‌دهند در حالی‌که رکوع‌کننده‌اند» — حالِ
  جمله‌ای: وَ هُمْ رَاكِعُونَ.}
\end{JKey}

\begin{JKey}[پاسخ‌نامهٔ تمرین‌ها]
\JAnswer{تمرینِ ١}{%
  ۱. اَلتِّقْنِيَّة ؛ ۲. اَلنَّفَق (جمع: اَلْأَنْفَاق) ؛ ۳. اَلتَّلّ (جمع:
  اَلتِّلَال) ؛ ۴. اَلْمَنْجَم (جمع: اَلْمَنَاجِم) ؛ ۵. اَلدَّؤُوب ؛
  ۶. اَلطُّنّ (جمع: اَلْأَطْنَان).}

\JAnswer{تمرینِ ٢}{%
  ترجمهٔ واژه‌های قرمز: أُعْطِيتُ = داده می‌شدم ؛ الْأَقَالِيمُ السَّبْعَةُ =
  سرزمین‌های هفتگانه ؛ أَفْلَاكِهَا = آسمان‌هایشان ؛ أَسْلُبُهَا = از آن
  (به زور) می‌گیرم ؛ جُلْبَ = پوستِ (گیاه/چوب) ؛ شَعِيرَةٍ = جویی ؛
  أَهْوَنُ = پست‌تر ؛ جَرَادَةٍ = ملخی ؛ تَقْضَمُهَا = آن را می‌جَوَد ؛
  نَعِيمٍ = نعمتی (مال/نعمت) ؛ سُبَاتِ الْعَقْلِ = به‌خواب‌رفتنِ خرد ؛
  قُبْحِ الزَّلَلِ = زشتیِ لغزش ؛ نَسْتَعِينُ = یاری می‌جوییم.\\[3pt]
  ۱. اعراب: نَمْلَةٍ: اسمِ مجرور به کسره (بعد از «فِي»؛ جارّ و مجرور) ؛
  جُلْبَ: مفعولِ «أَسْلُبُ»، منصوب به فتحه ؛ اَلْعَقْلِ: مضافٌ‌الیه، مجرور
  به کسره ؛ جَرَادَةٍ: مضافٌ‌الیه، مجرور به کسره.\\
  ۲. فَعَلْتُ: ماضی — صیغه: مفردِ متکلّم (مذکر)، معلوم.\\
  ۳. تَقْضَمُهَا: مضارع — صیغه: مفردِ مؤنثِ غایب، مرفوع (معلوم).\\
  ۴. دو جمعِ مکسّر: الْأَقَالِيم (مفرد: إِقْلِيم) و الْأَفْلَاك (مفرد:
  فَلَك).\\
  ۵. فعلِ مجهول: أُعْطِيتُ (داده شدم؛ بابِ إفعال، ماضیِ مجهول).\\
  ۶. تعدادِ فعل‌ها: ۹ فعل: أُعْطِيتُ ، أَعْصِيَ ، أَسْلُبُ ، فَعَلْتُ ،
  تَقْضَمُ ، يَفْنَى ، (لا) تَبْقَى ، نَعُودُ ، نَسْتَعِينُ.}

\JAnswer{تمرینِ ٣}{%
  واژهٔ ناهماهنگ در هر سطر: ۱. اَلْقَنَاة (توضیحِ سطر تعریفِ «اَلطِّين» است) ؛
  ۲. اَلْفِيزِيَاء (تعریفِ «حِزَامُ الْأَمَان» است) ؛ ۳. اَلطِّين (تعریفِ
  «اَلْفِيزِيَاء» است) ؛ ۴. اَلتِّلَال (تعریفِ «اَلْجَرَادَة» است) ؛
  ۵. حِزَامُ الْأَمَان (تعریفِ «اَلْقَنَاة» است) ؛ ۶. واژهٔ ششم فقط در
  فهرست است (تعریفِ ملخ در سطرِ چهارم آمد).}

\JAnswer{تمرینِ ٤}{همهٔ جمله‌ها با «آیا می‌دانی که...» آغاز می‌شوند:}
\JAnswer{ادامهٔ ٤}{%
  ۱. ... آیا می‌دانی که تلفظِ «گ» و «چ» و «ژ» در لهجه‌های عربیِ رایج بسیار
  وجود دارد؟\\
  ۲. ... که مغولان با وجودِ ساختنِ دیوارِ بزرگی پیرامونِ چین، توانستند بر آن
  هجوم آورند؟\\
  ۳. ... که نهنگ برای استخراجِ روغن از کبدش برای ساختِ موادّ آرایشی شکار
  می‌شود؟\\
  ۴. ... که خفاش تنها پستانداری است که می‌تواند پرواز کند؟\\
  ۵. ... که شمارِ مورچه‌های جهان تقریباً یک میلیون برابرِ انسان‌هاست؟\\
  ۶. ... که تیسفون، نزدیکِ بغداد، پایتختِ ساسانیان بود؟\\
  ۷. ... که اندازهٔ خرسِ پاندا هنگامِ تولد از موش کوچک‌تر است؟\\
  ۸. ... که زرافه گنگ است و تارهای صوتی ندارد؟\\
  ۹. ... که برگِ زیتون نمادِ صلح است؟}

\JAnswer{تمرینِ ٥}{گزینه‌های صحیح: ۱. (أ) ؛ ۲. (أ) ؛ ۳. (أ) ؛ ۴. (ب) ؛
۱ـ نَجَحَتْ: (ب) ؛ ۲ـ اَلطَّالِبَاتُ: (ب) ؛ ۳ـ اَلاِمْتِحَانِ: (أ).\\
توضیح: «عُمَّال» جمعِ تکسیرِ اسمِ فاعلِ «عَامِل» است و مبتداست؛
«الْمُجْتَهِدُونَ» اسمِ فاعلِ جمعِ مذکرِ سالمِ معرفه به «ال» و صفتِ
مرفوع؛ «يَشْتَغِلُونَ» فعلِ مضارعِ جمعِ مذکرِ غایب، ثلاثی مزیدِ بابِ
اِفتعال، معلوم و لازم و خبرِ مرفوع؛ «الْمَصْنَعِ» اسمِ مکانِ مفردِ معرفه و
مجرور به حرفِ جر؛ «نَجَحَتْ» ماضیِ مفردِ مؤنثِ غایب، ثلاثی مجرد، معلوم و
لازم؛ «الطَّالِبَاتُ» اسمِ فاعلِ جمعِ مؤنثِ سالمِ معرفه و فاعلِ مرفوع به
ضمّه؛ «الاِمْتِحَانِ» مصدرِ بابِ افتعال، معرفه به «ال» و مجرور به حرفِ جر.}

\JAnswer{تمرینِ ٦}{%
  ۱. تمرین را نوشت. / چرا درسِ (تکلیفِ) خودت را نمی‌نویسی؟ (خطابِ مؤنث) /
  چیزی ننوشتی.\\
  ۲. دو دوست (با هم) نامه‌نگاری کردند. / ای دو دوست، نامه‌نگاری کنید! /
  دو همکار نامه‌نگاری کردند.\\
  ۳. از موادّ قندی منع شدم. / ما را از بیرون رفتن بازندارید! / بازدارنده‌ای
  را در راه دیدیم.\\
  ۴. نگهبان از خواب خودداری کرده بود. / از مطالعه خودداری نکنید! / هرگز از
  بیرون رفتن خودداری نخواهیم کرد.\\
  ۵. چرا به تکلیف‌هایتان عمل نکردید؟ / آیا در کارخانه کار می‌کنید؟ /
  کارگران مشغولِ کارند.\\
  ۶. خدایا، ما را با فضلِ خودت رفتار کن. / خدایا، ما را با عدلِ خودت رفتار
  نکن. / آنان با ما خوب رفتار می‌کردند.\\
  ۷. مؤمن پروردگارش را یاد کرده است. / به نیکی یاد شدی. / استاد شاگردانِ
  قدیمی‌اش را یاد می‌کند.\\
  ۸. پدربزرگ و مادربزرگم مرا به یاد می‌آورند. / آموزگار به‌زودی ما را به یاد
  خواهد آورد. / تو را به یاد نمی‌آورم، ای هم‌کار (هم‌بازی) من.}

\JAnswer{تمرینِ ٧}{%
  ۱. خَاسِرًا (هر کس با دو چهره زندگی کند، زیان‌کار می‌میرد) ؛ ۲. مُقْتَدِرًا
  (تواناترین مردم کسی است که دشمنش را در حالِ توانایی می‌بخشد) ؛
  ۳. سَرِيعَةً (دشمنی هنگامِ بلاها به‌سرعت می‌رود) ؛ ۴. دو جملهٔ حالیّه با
  واو: وَ هُوَ يَضْحَكُ و وَ هُوَ يَبْكِي ؛ ۵. حَيًّا (نیکوکار زنده
  می‌ماند؛ جملهٔ «وَ إِنْ نُقِلَ...» نیز حالِ دوم است) ؛ ۶. وَحِيدًا (اگر
  می‌خواهی در کارت موفق شوی، تنها به آن برخیز و بر مردم توکّل مکن).}

\JAnswer{تمرینِ ٨}{%
  درست‌ها: ۳ ، ۴ ، ۸ ، ۹ ، ۱۰ ، ۱۱ ، ۱۳ ، ۱۴ ، ۱۵ ، ۱۷ ، ۱۸ ، ۱۹ ، ۲۰.\\
  نادرست‌ها: ۱ (صحیح: مَصَانِع) ؛ ۲ (صحیح: أَطْعِمَة؛ «مَطَاعِم» یعنی
  رستوران‌ها) ؛ ۵ (صحیح: عِظَام) ؛ ۶ (صحیح: مَمَالِك؛ «مُلُوك» جمعِ
  «مَلِيك» است) ؛ ۷ (صحیح: رُؤَسَاء؛ «رُؤُوس» جمعِ «رَأْس» است) ؛ ۱۲
  (صحیح: مَنَاجِم؛ «أَنْجُم» جمعِ «نَجْم» است) ؛ ۱۶ (صحیح: أَعْمَال؛
  «أَعْلَام» جمعِ «عَلَم» است).}
\end{JKey}
